Dans un modèle simple de l'effet de serre, l'atmosphère est une couche mince transparente au rayonnement solaire, mais qui absorbe tout le rayonnement infrarouge émis par le sol. Elle réémet ce rayonnement comme un corps noir, moitié vers l'espace, moitié vers le sol. La constante solaire est de \qty{\v{S}}{W/m^2} et l'albédo de la planète est de \num{\v{al}}.

\begin{enumerate}[label=\alph*)]
    \item En faisant le bilan d'énergie de la planète vue de l'espace, déterminez la température de la couche atmosphérique.
    \item En faisant le bilan d'énergie de la couche atmosphérique, montrez que $T_{\text{sol}} = 2^{1/4}\,T_{\text{atm}}$, puis calculez la température du sol.
    \item Quelle puissance par mètre carré le sol rayonne-t-il ? Comparez-la à la puissance solaire qu'il absorbe directement.
\end{enumerate}
